%%% BEGIN fig:method-coeff-minus-coll
\begin{figure}[htbp]
\centering
% photons, photon vertices and the cut, common to all graphs (origin: bottom centre of a graph)
\newcommand{\fbCMCphot}{%
  \draw[photon] (-1.22,1.95) -- (-0.88,1.4);
  \draw[photon] (0.88,1.4) -- (1.22,1.95);
  \node[vertex] at (-0.88,1.4) {}; \node[vertex] at (0.88,1.4) {};
  \draw[cut] (0,-0.18) -- (0,2.02);}
% collinear approximation: operator vertices, Wilson line between them, Born struck quark on top;
% #1 = position of the arrow on the Wilson line
\newcommand{\fbCMCapprox}[1]{%
  \draw[quark] (-0.88,0.92) -- (-0.88,1.4);
  \draw[quark=0.78] (-0.88,1.4) -- (0.88,1.4);
  \draw[quark] (0.88,1.4) -- (0.88,0.92);
  \draw[wilson=#1] (-0.88,0.92) -- (0.88,0.92);
  \node[opvertex] at (-0.88,0.92) {}; \node[opvertex] at (0.88,0.92) {};
  \fbCMCphot}
\begin{tikzpicture}[fd,
  fb lab/.style={font=\footnotesize,align=center,anchor=north,inner sep=1pt},
  fb row/.style={font=\small,anchor=east},
  fb frame/.style={rounded corners=3pt,line width=0.6pt}]
  % ======================= row 1: one-loop DIS cut graphs (gamma* q) =======================
  % (i) ISR squared
  \begin{scope}[xshift=0cm]
    \draw[quark] (-0.88,0) -- (-0.88,0.58);
    \draw[quark] (-0.88,0.58) -- (-0.88,1.4);
    \draw[quark=0.78] (-0.88,1.4) -- (0.88,1.4);
    \draw[quark] (0.88,1.4) -- (0.88,0.58);
    \draw[quark] (0.88,0.58) -- (0.88,0);
    \draw[gluon] (-0.88,0.58) -- (0.88,0.58);
    \fbCMCphot
    \node[vertex] at (-0.88,0.58) {}; \node[vertex] at (0.88,0.58) {};
    \node[font=\scriptsize,below,inner sep=1.5pt] at (-0.88,0) {$p$};
    \node[font=\scriptsize,right,inner sep=1.5pt] at (-1.12,1.98) {$q$};
    \node[fb lab] at (0,-0.36) {初态辐射\\的平方};
  \end{scope}
  % (ii) ISR--FSR interference
  \begin{scope}[xshift=2.9cm]
    \draw[quark] (-0.88,0) -- (-0.88,0.58);
    \draw[quark] (-0.88,0.58) -- (-0.88,1.4);
    \draw[quark=0.3] (-0.88,1.4) -- (0.46,1.4);
    \draw[quark] (0.46,1.4) -- (0.88,1.4);
    \draw[quark] (0.88,1.4) -- (0.88,0);
    \draw[gluon] (-0.88,0.58) to[out=0,in=-90] (0.46,1.4);
    \fbCMCphot
    \node[vertex] at (-0.88,0.58) {}; \node[vertex] at (0.46,1.4) {};
    \node[fb lab] at (0,-0.36) {ISR--FSR 干涉\\（+镜像）};
  \end{scope}
  % (iii) FSR squared
  \begin{scope}[xshift=5.8cm]
    \draw[quark] (-0.88,0) -- (-0.88,1.4);
    \draw[quark] (-0.88,1.4) -- (-0.46,1.4);
    \draw[quark=0.27] (-0.46,1.4) -- (0.46,1.4);
    \draw[quark] (0.46,1.4) -- (0.88,1.4);
    \draw[quark] (0.88,1.4) -- (0.88,0);
    \draw[gluon] (-0.46,1.4) arc[start angle=180,end angle=0,radius=0.46];
    \fbCMCphot
    \node[vertex] at (-0.46,1.4) {}; \node[vertex] at (0.46,1.4) {};
    \node[fb lab] at (0,-0.36) {末态辐射\\的平方};
  \end{scope}
  % (iv) virtual vertex correction
  \begin{scope}[xshift=8.7cm]
    \draw[quark] (-0.88,0) -- (-0.88,0.58);
    \draw[quark] (-0.88,0.58) -- (-0.88,1.4);
    \draw[quark] (-0.88,1.4) -- (-0.35,1.4);
    \draw[quark=0.65] (-0.35,1.4) -- (0.88,1.4);
    \draw[quark] (0.88,1.4) -- (0.88,0);
    \draw[gluon] (-0.88,0.58) to[out=0,in=-90] (-0.35,1.4);
    \fbCMCphot
    \node[vertex] at (-0.88,0.58) {}; \node[vertex] at (-0.35,1.4) {};
    \node[fb lab] at (0,-0.36) {顶点虚修正\\（+镜像）};
  \end{scope}
  % ======================= row 2: collinear approximations T_A Gamma = (PDF graph) x Born =======================
  \begin{scope}[yshift=-4.05cm]
    % (i') ladder graph of the PDF operator
    \begin{scope}[xshift=0cm]
      \draw[quark] (-0.88,0) -- (-0.88,0.42);
      \draw[quark] (-0.88,0.42) -- (-0.88,0.92);
      \draw[quark] (0.88,0.92) -- (0.88,0.42);
      \draw[quark] (0.88,0.42) -- (0.88,0);
      \draw[gluon] (-0.88,0.42) -- (0.88,0.42);
      \fbCMCapprox{0.25}
      \node[vertex] at (-0.88,0.42) {}; \node[vertex] at (0.88,0.42) {};
      \node[font=\scriptsize,below,inner sep=1.5pt] at (-0.88,0) {$p$};
      \node[fb lab] at (0,-0.36) {PDF 算符的\\梯形图};
    \end{scope}
    % (ii') real Wilson-line graph
    \begin{scope}[xshift=2.9cm]
      \draw[quark] (-0.88,0) -- (-0.88,0.42);
      \draw[quark] (-0.88,0.42) -- (-0.88,0.92);
      \draw[quark] (0.88,0.92) -- (0.88,0);
      \draw[gluon] (-0.88,0.42) to[out=0,in=-90] (0.46,0.92);
      \fbCMCapprox{0.22}
      \node[vertex] at (-0.88,0.42) {}; \node[vertex] at (0.46,0.92) {};
      \node[fb lab] at (0,-0.36) {Wilson 线实图\\（+镜像）};
    \end{scope}
    % (iii') no target-collinear region
    \begin{scope}[xshift=5.8cm]
      \node[font=\large] at (0,0.9) {$0$};
      \node[fb lab] at (0,-0.36) {无靶共线区};
    \end{scope}
    % (iv') virtual Wilson-line graph
    \begin{scope}[xshift=8.7cm]
      \draw[quark] (-0.88,0) -- (-0.88,0.42);
      \draw[quark] (-0.88,0.42) -- (-0.88,0.92);
      \draw[quark] (0.88,0.92) -- (0.88,0);
      \draw[gluon] (-0.88,0.42) to[out=0,in=-90] (-0.35,0.92);
      \fbCMCapprox{0.72}
      \node[vertex] at (-0.88,0.42) {}; \node[vertex] at (-0.35,0.92) {};
      \node[fb lab] at (0,-0.36) {Wilson 线虚图\\（+镜像）};
    \end{scope}
  \end{scope}
  % ======================= colour frames, row labels, result =======================
  \begin{scope}[on background layer]
    \foreach \x/\c in {0/green,2.9/blue,5.8/gray,8.7/orange}
      \draw[fb frame,draw=\c!55!black,fill=\c!6] (\x-1.4,-5.3) rectangle (\x+1.4,2.2);
  \end{scope}
  \node[fb row] at (-1.5,0.9) {$\Gamma^{(1)}$};
  \node[fb row] at (-1.5,-3.15) {$-\,T_A\Gamma^{(1)}$};
  \foreach \x in {0,2.9,5.8,8.7} \node[font=\large] at (\x,-1.62) {$-$};
  \draw[decorate,decoration={brace,amplitude=5pt},line width=0.6pt] (10.25,2.15) -- (10.25,-5.25);
  \node[anchor=west,align=left,font=\small] at (10.45,-1.55) {各列之和\\[3pt]$=C_H\Gamma^{(1)}$\\[3pt]$=\dfrac{\as}{2\pi}\,C_2^q$\\[1pt]{\footnotesize（$\mu=Q$）}};
\end{tikzpicture}
\caption{单圈 DIS 系数函数 = 割图减去其共线近似（$\gamma^*q$ 道，Feynman 规范；竖直红色虚线为末态割线，左侧为振幅）。
上行为 $\order{\as}$ 割图 $\Gamma^{(1)}$，下行为其靶共线近似 $T_A\Gamma^{(1)}=f^{(1)}_{q/q}\otimes{\rm Born}$：PDF 算符的割图（$\otimes$ 为算符顶点，双线为沿 $n_B$ 的 Wilson 线）乘以领头阶硬部分；同色框中的图一一对应。
初态辐射的平方对应 PDF 算符的梯形图，ISR--FSR 干涉对应胶子连到割线另一侧 Wilson 线的实图，顶点虚修正对应同侧的 Wilson 线虚图（外线自能无标度，未画出）；末态辐射没有靶共线区，不需减除。
各列之差的和即硬区减除项 $C_H\Gamma^{(1)}$，在 $\mu=Q$ 时为 $\frac{\as}{2\pi}C_2^q$；纯维数正规化中下行各图无标度，只留下 $\MSbar$ 抵消项，$C_A\Gamma^{(1)}=-\frac{\as}{2\pi}\frac1\varepsilon P_{qq}$（VF-DIS-6，VF-PDF-3）。}
\label{fig:method-coeff-minus-coll}
\end{figure}
%%% END fig:method-coeff-minus-coll
